\section{The optimized Ricci-box obstruction}
\label{sec:ricci-box-optimization}
Let the actual monopole supply give a positive budget $A$ and the mixed
curvature inequality
\begin{equation}\label{eq:box-input}
 \int_M(s-\sqrt6|W^+|)^2\ge72\pi^2A.
\end{equation}
Write $T=(1-\eta)A$. The symbol $\eta$ denotes a topological ratio;
it is not the projected class square $p$. Suppose
\[
 d=\|g^{-1}\Ric+I\|_{L^\infty,\mathrm{op}}<\tfrac12.
\]
The four Ricci eigenvalues are $-1+a_i$, with $|a_i|\le d$. Put
$w=|W^+|$ and $F=s^2/24-|\Ric^\circ|^2/2$. Then
\begin{equation}\label{eq:box-cw}
 \int_M(2w^2+F)=4\pi^2(1-\eta)A.
\end{equation}
Keeping the correlation between scalar and trace-free Ricci curvature
improves the estimate obtained by maximizing them separately.

\begin{proposition}\label{prop:box-optimization}
The exact algebraic supremum over the Ricci box and $w\ge0$ is
\begin{equation}\label{eq:box-constant}
 \sup\frac{(s-\sqrt6w)^2}{2w^2+F}=C(d)=
 \begin{cases}
 \displaystyle\frac{9(3-d^2)}{1-3d^2},&0\le d\le2-\sqrt3,\\[2mm]
 \displaystyle\frac{27(1-d)}{(1-2d)(1+d)},&2-\sqrt3\le d<\tfrac12.
 \end{cases}
\end{equation}
All the denominators in this expression are positive.
\end{proposition}
\begin{proof}
Set $x_i=1-a_i>0$ and $S_0=\sum_i x_i=-s$. The ratio
\[
 R_0=\frac{|\Ric^\circ|^2}{s^2}
     =\frac{\sum_i x_i^2}{(\sum_i x_i)^2}-\frac14
\]
attains its maximum at a vertex of $[1-d,1+d]^4$. Indeed, fixing the
other coordinates with sum $u>0$ and square sum $v>0$, the derivative
of $(x^2+v)/(x+u)^2$ is $2(ux-v)/(x+u)^3$; every interior critical
point is a minimum. Moving one coordinate at a time to an endpoint
therefore cannot decrease a maximum. If $m$ coordinates equal $1+d$,
the vertex value is
\[
 R_m=\frac{m(4-m)d^2}{[4+(2m-4)d]^2}.
\]
The values for $m=0,4$ vanish and $R_3\le R_1$. Consequently
\begin{equation}\label{eq:box-ratio}
 R_0\le R(d)=\max\left\{\frac{d^2}{4},
                     \frac{3d^2}{4(2-d)^2}\right\}<\frac1{12}.
\end{equation}
The last inequality follows directly from $d<1/2$; hence
$F=s^2(1-12R_0)/24>0$. For fixed Ricci eigenvalues, weighted
Cauchy--Schwarz gives
\[
 \frac{(S_0+\sqrt6w)^2}{F+2w^2}
 \le\frac{S_0^2}{F}+3=\frac{24}{1-12R_0}+3.
\]
Equality occurs at $w=\sqrt6F/(2S_0)$. The two maxima in
\eqref{eq:box-ratio} coincide at $d=2-\sqrt3$; substituting them
gives \eqref{eq:box-constant}. They correspond respectively to two
large and two small eigenvalue magnitudes, and to one large and three
small magnitudes.
\end{proof}

\begin{corollary}\label{cor:box-topological}
If $1/3<\eta<1$, every metric on a smooth type with the stated
all-metric monopole supply satisfies $\delta_-(g)\ge d_*(\eta)$, where
\begin{equation}\label{eq:box-inverse}
 d_*(\eta)=
 \begin{cases}
 \displaystyle\sqrt{\frac{3\eta-1}{\eta+5}},
 &\frac13<\eta\le2\sqrt3-3,\\[2mm]
 \displaystyle\frac{\sqrt{(3\eta-1)(3\eta+15)}-(3\eta-1)}8,
 &2\sqrt3-3\le\eta<1.
 \end{cases}
\end{equation}
For $\eta\ge1$, the corresponding bound is $\delta_-(g)\ge1/2$.
\end{corollary}
\begin{proof}
For $d<1/2$, combine \eqref{eq:box-input}--\eqref{eq:box-constant}:
$1-\eta\ge18/C(d)$. Equivalently,
\[
 \eta\le H(d):=
 \begin{cases}
 (1+5d^2)/(3-d^2),&d\le2-\sqrt3,\\
 \frac13+4d^2/[3(1-d)],&d\ge2-\sqrt3.
 \end{cases}
\]
Both branches are increasing, their common value at the junction is
$2\sqrt3-3$, and their inversion gives \eqref{eq:box-inverse}.
Since $d_*(\eta)<1/2$, metrics with $d\ge1/2$ satisfy the same bound.
Finally $F>0$ makes the left side of \eqref{eq:box-cw} positive,
excluding $d<1/2$ whenever $\eta\ge1$.
\end{proof}

For comparison, the separated estimate gives
$H_{\rm sep}(d)=1/3+2d^2/(1-d)^2$. Direct subtraction yields
\[
 H_{\rm sep}(d)-H(d)=
 \begin{cases}
 \displaystyle\frac{2d^2(1+16d-11d^2)}{3(3-d^2)(1-d)^2},
       &0<d\le2-\sqrt3,\\[2mm]
 \displaystyle\frac{2d^2(1+2d)}{3(1-d)^2},
       &2-\sqrt3\le d<\tfrac12.
 \end{cases}
\]
Thus the improvement is strict throughout the nonzero defect range.
As $q=\eta-1/3\downarrow0$, the new lower bound is
$d_*(\eta)=\frac34\sqrt q+O(q^{3/2})$. This calculation is sharp for
the algebraic quotient, not a construction of closed equality metrics.
It covers ratios above $1/3$ and does not supply the uniform differential
separation at or below that threshold proved by the main calibration
argument.
